\documentclass[12pt,a4paper]{article}
\usepackage[utf8]{inputenc}
\usepackage[T1]{fontenc}
\usepackage[french]{babel}
\usepackage{geometry}
\usepackage{xcolor}
\usepackage{fancyhdr}

\geometry{margin=2cm}
\pagestyle{fancy}
\fancyhf{}
\fancyhead[L]{CEJM - Fiche de Révision}
\fancyhead[R]{Chapitre 4}
\fancyfoot[C]{\thepage}
\setlength{\headheight}{15pt}

% Couleurs
\definecolor{bleu}{RGB}{0,90,170}
\definecolor{vert}{RGB}{0,150,0}
\definecolor{rouge}{RGB}{200,0,0}
\definecolor{gris}{RGB}{240,240,240}

% Commandes pour encadrés
\newcommand{\definition}[2]{\noindent\fcolorbox{bleu}{gris}{\parbox{\dimexpr\textwidth-2\fboxsep-2\fboxrule\relax}{\textbf{\textcolor{bleu}{DÉFINITION :}} #1\\\textit{#2}}}}

\newcommand{\pointcle}[1]{\noindent\fcolorbox{vert}{white}{\parbox{\dimexpr\textwidth-2\fboxsep-2\fboxrule\relax}{\textbf{\textcolor{vert}{POINT CLÉ :}} #1}}}

\newcommand{\exemple}[1]{\noindent\fcolorbox{rouge}{white}{\parbox{\dimexpr\textwidth-2\fboxsep-2\fboxrule\relax}{\textbf{\textcolor{rouge}{EXEMPLE :}} #1}}}

\title{\textbf{Fiche de Révision CEJM}\\[2mm]\large Chapitre 4 : La création d'entreprise et sa performance}
\author{}
\date{}

\begin{document}
\maketitle

\section*{Problématique}
\textbf{Comment créer une entreprise et mesurer sa performance économique et financière ?}

\section{Définitions essentielles}

\definition{Entreprise}{Organisation économique qui combine des facteurs de production pour produire des biens ou services destinés à la vente.}

\vspace{3mm}
\definition{Entrepreneur}{Personne qui crée et dirige une entreprise, prend des risques économiques en vue de réaliser un profit.}

\vspace{3mm}
\definition{Business plan}{Document de synthèse qui présente le projet d'entreprise, sa stratégie et ses prévisions financières.}

\vspace{3mm}
\definition{Chiffre d'affaires}{Montant total des ventes réalisées par l'entreprise sur une période donnée.}

\vspace{3mm}
\definition{Performance}{Capacité de l'entreprise à atteindre ses objectifs de manière efficace et efficiente.}

\section{Le processus de création d'entreprise}

\subsection{L'idée et l'opportunité}
\pointcle{Identifier une opportunité de marché et valider l'idée}
\begin{itemize}
    \item \textbf{Sources d'idées} : expérience, innovation, imitation, franchise
    \item \textbf{Étude de marché} : analyser la demande, la concurrence, l'environnement
    \item \textbf{Validation} : faisabilité technique, commerciale, financière
\end{itemize}

\subsection{Le business plan}
\pointcle{Document stratégique indispensable pour convaincre les financeurs}
\begin{itemize}
    \item \textbf{Présentation du projet} : produit/service, marché, équipe
    \item \textbf{Stratégie commerciale} : positionnement, marketing-mix
    \item \textbf{Prévisions financières} : compte de résultat, plan de financement
    \item \textbf{Planning} : étapes de réalisation, calendrier
\end{itemize}

\subsection{Le financement}
\pointcle{Réunir les capitaux nécessaires au lancement}
\begin{itemize}
    \item \textbf{Apports personnels} : épargne, famille, amis
    \item \textbf{Aides publiques} : subventions, ACRE, prêts d'honneur
    \item \textbf{Emprunts bancaires} : crédits professionnels
    \item \textbf{Investisseurs} : business angels, capital-risque
\end{itemize}

\section{Les formes juridiques d'entreprise}

\subsection{Entreprise individuelle}
\pointcle{Forme la plus simple, patrimoine personnel et professionnel confondus}
\begin{itemize}
    \item \textbf{Avantages} : simplicité, rapidité, coûts réduits
    \item \textbf{Inconvénients} : responsabilité illimitée, difficulté de financement
    \item \textbf{Variantes} : micro-entreprise, EIRL
\end{itemize}

\subsection{Sociétés}
\pointcle{Personnalité juridique distincte, responsabilité limitée}

\subsubsection{SARL (Société à Responsabilité Limitée)}
\begin{itemize}
    \item 1 à 100 associés
    \item Capital minimum : 1 €
    \item Responsabilité limitée aux apports
    \item Dirigeant : gérant
\end{itemize}

\subsubsection{SAS (Société par Actions Simplifiée)}
\begin{itemize}
    \item 1 associé minimum (SASU si seul)
    \item Grande flexibilité statutaire
    \item Dirigeant : président
    \item Facilite l'entrée d'investisseurs
\end{itemize}

\subsubsection{SA (Société Anonyme)}
\begin{itemize}
    \item Minimum 2 associés (7 si cotée)
    \item Capital minimum : 37 000 €
    \item Pour les grandes entreprises
    \item Dirigeant : directeur général
\end{itemize}

\section{La mesure de la performance}

\subsection{Performance commerciale}
\pointcle{Capacité à vendre et à fidéliser la clientèle}
\begin{itemize}
    \item \textbf{Chiffre d'affaires} : CA = Prix × Quantités vendues
    \item \textbf{Parts de marché} : CA entreprise / CA total du marché
    \item \textbf{Taux de croissance} : (CA n - CA n-1) / CA n-1 × 100
    \item \textbf{Panier moyen} : CA / Nombre de clients
\end{itemize}

\subsection{Performance financière}
\pointcle{Capacité à générer des profits et à gérer les finances}

\subsubsection{Résultats}
\begin{itemize}
    \item \textbf{Résultat d'exploitation} : Produits d'exploitation - Charges d'exploitation
    \item \textbf{Résultat net} : Résultat après impôts et charges financières
    \item \textbf{Marge commerciale} : Ventes - Coût d'achat des marchandises vendues
\end{itemize}

\subsubsection{Ratios de rentabilité}
\begin{itemize}
    \item \textbf{Taux de marge} : Résultat / CA × 100
    \item \textbf{Rentabilité économique} : Résultat d'exploitation / Capitaux investis
    \item \textbf{Rentabilité financière} : Résultat net / Capitaux propres
\end{itemize}

\subsection{Performance sociale}
\pointcle{Capacité à gérer les ressources humaines}
\begin{itemize}
    \item \textbf{Productivité} : Production / Nombre de salariés
    \item \textbf{Taux d'absentéisme} : Heures d'absence / Heures théoriques
    \item \textbf{Turn-over} : Nombre de départs / Effectif moyen
    \item \textbf{Formation} : Budget formation / Masse salariale
\end{itemize}

\section{Les facteurs de performance}

\subsection{Facteurs internes}
\begin{itemize}
    \item \textbf{Innovation} : nouveaux produits, nouvelles technologies
    \item \textbf{Qualité} : satisfaction client, certification
    \item \textbf{Organisation} : structure efficace, processus optimisés
    \item \textbf{Ressources humaines} : compétences, motivation
\end{itemize}

\subsection{Facteurs externes}
\begin{itemize}
    \item \textbf{Marché} : taille, croissance, concurrence
    \item \textbf{Environnement économique} : conjoncture, taux d'intérêt
    \item \textbf{Réglementation} : contraintes légales, normes
    \item \textbf{Technologie} : évolutions, disruptions
\end{itemize}

\section{Exemples concrets}

\exemple{Start-up technologique : financement par capital-risque, croissance rapide du CA, SAS pour attirer investisseurs}

\exemple{Commerce de proximité : faible capital de départ, performance mesurée par CA journalier et panier moyen}

\exemple{PME industrielle : investissements lourds, performance liée à la productivité et aux marges}

\section{Les difficultés de l'entreprise}

\subsection{Causes de défaillance}
\begin{itemize}
    \item \textbf{Insuffisance de fonds propres} : sous-capitalisation
    \item \textbf{Mauvaise gestion} : pilotage défaillant
    \item \textbf{Marché inadapté} : mauvaise étude de marché
    \item \textbf{Concurrence} : pression concurrentielle
\end{itemize}

\subsection{Signaux d'alerte}
\begin{itemize}
    \item Baisse du CA
    \item Dégradation des marges
    \item Difficultés de trésorerie
    \item Retards de paiement
\end{itemize}

\section{Points clés à retenir}

\pointcle{La création d'entreprise nécessite une préparation rigoureuse et un business plan solide}

\pointcle{Le choix de la forme juridique dépend du projet, des associés et du financement}

\pointcle{La performance se mesure sur plusieurs dimensions : commerciale, financière, sociale}

\pointcle{L'innovation et l'adaptation au marché sont clés pour la performance}

\pointcle{Le pilotage par les indicateurs permet d'anticiper les difficultés}

\section{Mots-clés à connaître}
\begin{itemize}
    \item Entrepreneur/Entrepreneuriat
    \item Business plan
    \item Forme juridique
    \item SARL/SAS/SA
    \item Chiffre d'affaires
    \item Résultat net
    \item Rentabilité
    \item Part de marché
    \item Productivité
    \item Innovation
    \item Performance
    \item Capital social
\end{itemize}

\end{document}